\section{Mid-Term Examination --- Set A (30 Marks)}
\label{sec:midterm_set_a}

\LPUPaperHeader{Mid-Term Examination}{A}{90 Minutes}{30 Marks}{Units 1, 2, and 3 (Matrices and ODEs)}

\begin{center}
\sffamily\textbf{\color{AlertRed}Note: 30 Multiple Choice Questions. Each question carries 1 mark. Negative marking $-0.25$ marks.}
\end{center}

\begin{enumerate}[label=\textbf{\arabic*.}]
    \item The rank of an $n \times n$ non-singular matrix is:
    \begin{enumerate*}[label=(\alph*)]
        \item $n$ \quad
        \item $0$ \quad
        \item $n - 1$ \quad
        \item $1$
    \end{enumerate*}
    \item The sum of eigenvalues of a matrix equals its:
    \begin{enumerate*}[label=(\alph*)]
        \item Determinant \quad
        \item Trace \quad
        \item Rank \quad
        \item Nullity
    \end{enumerate*}
    \item If the Wronskian of two solutions is zero everywhere on an interval, the solutions are:
    \begin{enumerate*}[label=(\alph*)]
        \item Linearly independent \quad
        \item Linearly dependent \quad
        \item Orthogonal \quad
        \item Singular
    \end{enumerate*}
    \item The auxiliary equation for $\frac{d^2 y}{dx^2} + 9y = 0$ is:
    \begin{enumerate*}[label=(\alph*)]
        \item $m^2 - 9 = 0$ \quad
        \item $m^2 + 9 = 0$ \quad
        \item $m + 9 = 0$ \quad
        \item $m^2 + 3 = 0$
    \end{enumerate*}
    \item The particular integral for $(D - 2)y = e^{2x}$ is:
    \begin{enumerate*}[label=(\alph*)]
        \item $e^{2x}$ \quad
        \item $x e^{2x}$ \quad
        \item $\frac{1}{2}e^{2x}$ \quad
        \item $x^2 e^{2x}$
    \end{enumerate*}
    \item For the Cauchy-Euler equation $x^2 y'' + 2xy' - 6y = 0$, the substitution used is:
    \begin{enumerate*}[label=(\alph*)]
        \item $x = e^z$ \quad
        \item $x = \sin z$ \quad
        \item $y = e^z$ \quad
        \item $x = z^2$
    \end{enumerate*}
\end{enumerate}

\begin{center}
\textbf{[Questions 7 to 30 continue across Units 1 to 3 covering all core theorems and methods]}
\end{center}

\subsection*{Mid-Term Answer Key \& Detailed Solutions}
\begin{answerkeytable}{Mid-Term Set A Answer Key}
\begin{tabular}{|c|c|c|c|c|c|}
\hline
\textbf{Q1} & \textbf{Q2} & \textbf{Q3} & \textbf{Q4} & \textbf{Q5} & \textbf{Q6} \\
\hline
(a) & (b) & (b) & (b) & (b) & (a) \\
\hline
\end{tabular}
\end{answerkeytable}

\begin{solutionbox}[Explanations for Core Questions]
$\bullet$ \textbf{Q1}: A non-singular $n \times n$ matrix has $\det(A) \neq 0$, meaning all $n$ rows are linearly independent; hence rank is $n$.\\
$\bullet$ \textbf{Q2}: By Vieta's formulas on the characteristic polynomial, the coefficient of $\lambda^{n-1}$ is $(-1)\sum a_{ii} = (-1)\operatorname{tr}(A)$, hence $\sum \lambda_i = \operatorname{tr}(A)$.\\
$\bullet$ \textbf{Q5}: Since $D = 2$ is a root of the auxiliary operator $f(D) = D - 2$, the standard formula fails; applying the failure rule: $y_p = x \frac{1}{f'(D)}e^{2x} = x \frac{1}{1}e^{2x} = x e^{2x}$.
\end{solutionbox}

\newpage
